% TOP_PROBLEM: {"id": 2800704, "title": "10 Lectures and 42 Open Problems — The Deletion Channel", "category_id": 15, "category": {"id": 15, "name": "computer_science", "display_name": "Computer Science", "description": "Computational complexity, algorithms, and theoretical CS.", "slug": "computer-science", "order_index": 15, "created_at": "2026-07-31T15:26:25.671Z"}, "status": "partially_solved", "problem_number": "AMR-027-0704", "set_id": 13, "statusQualification": "The source says to draw a random binary input. Average-case and worst-case trace complexity, and the chosen error probabilities, are explicitly distinguished."}
% FIRST_POSTED: 2026-10-06
% ATTEMPT_STATUS: unresolved
% OPENAI_CLAIM: {"schema_version":1,"source_repo":"https://github.com/openai/math","source_commit":"adc7f1241b42e322a6451854ab7e4b4c146bf78a","release_date":"2026-10-06","match":"partial","resolution":"partial","independently_reviewed":false,"second_pass":"n/a","families":[{"family":"122","title":"Quantitative trace-reconstruction bounds with a uniform decoder","subject":"Theoretical computer science","manuscripts":[{"title":"A latest-anchor induction with spectrally compact masks for worst-case trace reconstruction","pdf_path":"preprints/A-latest-anchor-induction-with-spectrally-compact-masks-for-worst-case-trace-reconstruction-October-5-2026/paper.pdf","readme_path":"preprints/A-latest-anchor-induction-with-spectrally-compact-masks-for-worst-case-trace-reconstruction-October-5-2026/README.md","date":"2026-10-05","theorem_ref":"Theorem 1"},{"title":"Uniform quasipolynomial-time trace reconstruction","pdf_path":"preprints/Uniform-quasipolynomial-time-trace-reconstruction-October-5-2026/uniform-trace-reconstruction.pdf","readme_path":"preprints/Uniform-quasipolynomial-time-trace-reconstruction-October-5-2026/README.md","date":"2026-10-05","theorem_ref":"Theorem 1.1"},{"title":"Quantitative lower bounds for trace reconstruction","pdf_path":"preprints/quantitative-lower-bounds-for-trace-reconstruction-September-24-2026/paper.pdf","readme_path":"preprints/quantitative-lower-bounds-for-trace-reconstruction-September-24-2026/README.md","date":"2026-09-24","theorem_ref":"Theorems 1.1 and 1.2"}],"lean":{"doc_path":"lean/docs/122.md","comparators":["lean/ComparatorChallenges/TraceReconstruction.json","lean/ComparatorChallenges/TraceReconstruction.lean"],"declarations":[],"covers_main_theorem":false},"reasoning_trace":null}]}
\documentclass[11pt]{article}
\usepackage[margin=25mm]{geometry}
\usepackage{fontspec}
\setmainfont{Latin Modern Roman}
\usepackage{amsmath,amssymb,mathtools}
\usepackage{unicode-math}
\setmathfont{Latin Modern Math}
\usepackage{enumitem}
\usepackage{xurl,hyperref,bookmark}
\hypersetup{colorlinks=true,urlcolor=blue}
\setcounter{secnumdepth}{0}
\setlength{\parindent}{0pt}
\setlength{\parskip}{5pt}
\setlength{\emergencystretch}{3em}
\begin{document}
\section*{10 Lectures and 42 Open Problems — The Deletion Channel}

\subsection{OpenAI's claim}
OpenAI presents the following claims relevant to this question, with manuscript numbers following paper order in Sources. In manuscript 1 (family 122; 2026-10-05), the asserted result is that worst-case trace reconstruction has quasipolynomial sample bounds with logarithm $O((\log n)^3(1+\log\log(2n))^6)$ at fixed retention, including general exactly observed symbols, without an algorithmic guarantee; In manuscript 2 (family 122; 2026-10-05), the asserted result is that every binary word can be reconstructed with known rational retention probability using a uniform algorithm with quasipolynomial sample and bit bounds at fixed retention, and polynomial bounds for sufficiently small deletion; In manuscript 3 (family 122; 2026-09-24), the asserted result is that worst-case exact binary trace reconstruction needs $n^{\Omega(\log\log n)}$ samples at any fixed deletion probability, even with unrestricted computation.

\subsection{Outline of the argument}
For manuscript 1 (family 122; 2026-10-05), gap probes use compact spectral masks and a latest-anchor induction to keep a signal during shuffling and scale transfer. Radial compactification and stopping-stage interpolation produce separation of trace laws; the argument then reduces general symbols to the detected binary distinction.

For manuscript 2 (family 122; 2026-10-05), low-order trace moments propagate a distinguishing signal through scales using damped products and prescribed frequencies. Positive moment constraints recover prefixes, and a bounded feasibility iteration turns the statistical separation into an explicit finite-precision decoder.

For manuscript 3 (family 122; 2026-09-24), nested signed words hide fair orientation bits inside overlapping windows. An auxiliary observation retains side data while projecting to genuine traces; likelihood-score induction forces a large binary subtree for every surviving term and bounds the distance between the two trace laws.

\subsection{Scope relative to this problem}
These results concern worst-case input words and hard-pair separation. They do not determine the typical random-word complexity or the exact extremal pair growth, so those additional clauses of the deletion-channel problem remain unanswered.

\subsection{Formal verification}
For manuscript 3 (family 122; 2026-09-24), OpenAI lists Lean statements for this assertion: worst-case exact binary trace reconstruction needs $n^{\Omega(\log\log n)}$ samples at any fixed deletion probability, even with unrestricted computation.

The selected statement includes the clause used for this problem.

No Lean formalization is listed at this snapshot for manuscript 1 (family 122; 2026-10-05), manuscript 2 (family 122; 2026-10-05).

The linked family Comparator checks are \texttt{TraceReconstruction.lean}.

Their named declarations include \texttt{OAI.TraceReconstruction.quantitative\_sample\_lower\_bound}, \texttt{OAI.TraceReconstruction.superpolynomial\_one\_trace}, \texttt{OAI.TraceReconstruction.superpolynomial\_sample\_complexity}.

This site has not executed the Lean code. Described statement coverage is not an independent proof check.

\subsection{Status caveat}
The manuscripts were produced by an internal OpenAI model. OpenAI cautions that results without formalization can contain mistakes. This site has not independently checked the proofs; the label records the scope of OpenAI's claim.

\subsection{Sources}
\begin{itemize}
\item OpenAI, \emph{A latest-anchor induction with spectrally compact masks for worst-case trace reconstruction}, 2026-10-05; citation key \texttt{OAI:A-latest-anchor-induction-with-spectrally-compact-masks-for-worst-case-trace-reconstruction-October-5-2026}. \href{https://github.com/openai/math/blob/adc7f1241b42e322a6451854ab7e4b4c146bf78a/preprints/A-latest-anchor-induction-with-spectrally-compact-masks-for-worst-case-trace-reconstruction-October-5-2026/paper.pdf}{PDF (GitHub)}; \href{https://github.com/openai/math/raw/adc7f1241b42e322a6451854ab7e4b4c146bf78a/preprints/A-latest-anchor-induction-with-spectrally-compact-masks-for-worst-case-trace-reconstruction-October-5-2026/paper.pdf}{PDF (direct)}; \href{https://github.com/openai/math/blob/adc7f1241b42e322a6451854ab7e4b4c146bf78a/preprints/A-latest-anchor-induction-with-spectrally-compact-masks-for-worst-case-trace-reconstruction-October-5-2026/README.md}{Manuscript page and citation}; \href{https://github.com/openai/math/blob/main/preprints/A-latest-anchor-induction-with-spectrally-compact-masks-for-worst-case-trace-reconstruction-October-5-2026/paper.pdf}{Latest version}.
\item OpenAI, \emph{Uniform quasipolynomial-time trace reconstruction}, 2026-10-05; citation key \texttt{OAI:Uniform-quasipolynomial-time-trace-reconstruction-October-5-2026}. \href{https://github.com/openai/math/blob/adc7f1241b42e322a6451854ab7e4b4c146bf78a/preprints/Uniform-quasipolynomial-time-trace-reconstruction-October-5-2026/uniform-trace-reconstruction.pdf}{PDF (GitHub)}; \href{https://github.com/openai/math/raw/adc7f1241b42e322a6451854ab7e4b4c146bf78a/preprints/Uniform-quasipolynomial-time-trace-reconstruction-October-5-2026/uniform-trace-reconstruction.pdf}{PDF (direct)}; \href{https://github.com/openai/math/blob/adc7f1241b42e322a6451854ab7e4b4c146bf78a/preprints/Uniform-quasipolynomial-time-trace-reconstruction-October-5-2026/README.md}{Manuscript page and citation}; \href{https://github.com/openai/math/blob/main/preprints/Uniform-quasipolynomial-time-trace-reconstruction-October-5-2026/uniform-trace-reconstruction.pdf}{Latest version}.
\item OpenAI, \emph{Quantitative lower bounds for trace reconstruction}, 2026-09-24; citation key \texttt{OAI:quantitative-lower-bounds-for-trace-reconstruction-September-24-2026}. \href{https://github.com/openai/math/blob/adc7f1241b42e322a6451854ab7e4b4c146bf78a/preprints/quantitative-lower-bounds-for-trace-reconstruction-September-24-2026/paper.pdf}{PDF (GitHub)}; \href{https://github.com/openai/math/raw/adc7f1241b42e322a6451854ab7e4b4c146bf78a/preprints/quantitative-lower-bounds-for-trace-reconstruction-September-24-2026/paper.pdf}{PDF (direct)}; \href{https://github.com/openai/math/blob/adc7f1241b42e322a6451854ab7e4b4c146bf78a/preprints/quantitative-lower-bounds-for-trace-reconstruction-September-24-2026/README.md}{Manuscript page and citation}; \href{https://github.com/openai/math/blob/main/preprints/quantitative-lower-bounds-for-trace-reconstruction-September-24-2026/paper.pdf}{Latest version}.
\item Family 122: \href{https://github.com/openai/math/blob/adc7f1241b42e322a6451854ab7e4b4c146bf78a/lean/docs/122.md}{Lean scope}.
\item \href{https://github.com/openai/math/blob/adc7f1241b42e322a6451854ab7e4b4c146bf78a/lean/ComparatorChallenges/TraceReconstruction.json}{Comparator TraceReconstruction.json}.
\item \href{https://github.com/openai/math/blob/adc7f1241b42e322a6451854ab7e4b4c146bf78a/lean/ComparatorChallenges/TraceReconstruction.lean}{Comparator TraceReconstruction.lean}.
\item \href{https://github.com/openai/math/blob/adc7f1241b42e322a6451854ab7e4b4c146bf78a/CONTENTS.md}{OpenAI catalogue entry}.
\end{itemize}
\end{document}
